\section{Results}
\label{sec:results}

% OWNER: Sanjana Garimella (results/discussion per proposal)
% Figures reference live files in ../results/ via \graphicspath in main.tex.

We evaluate every model on the same leak-free set of $19{,}819$ minimal pairs
($1{,}528$ with an attractor, $18{,}291$ without), held identical across all
five noise conditions. Unless noted, reported values are the mean over six LSTM
runs per condition (seeds $\{1,2,3\}$, run on both local MPS and Colab T4);
$\pm$ denotes standard deviation across runs. The RNNG arm
(Section~\ref{sec:rnng}) is a single full-scale seed evaluated on the same
held-out set, with additional seeds for error bars still in progress.

\subsection{Noise degrades agreement, with a threshold near 50\%}

Overall minimal-pair accuracy is remarkably stable across the first four noise
levels and then collapses at the highest rate (Table~\ref{tab:main},
Figure~\ref{fig:acc_vs_noise}). Accuracy falls only $2.7$ points as the error
rate climbs from $0\%$ to $10\%$ ($0.961 \to 0.934$), but drops $22.6$ points
between $10\%$ and $50\%$ ($0.934 \to 0.708$). The model thus tolerates
substantial agreement noise, up to roughly one corrupted verb in ten, before
its grasp of the rule breaks down. This non-linear, threshold-like profile
matches the first outcome anticipated in our proposal: input quality can be
degraded considerably before grammar learning is affected, but beyond a critical
point performance falls sharply toward chance.

\begin{table}[t]
  \centering
  \small
  \begin{tabular}{lccc}
    \toprule
    Condition & Rate & Acc.\ (all) & Attractor gap \\
    \midrule
    baseline     & 0.000 & $0.961 \pm 0.006$ & $0.073 \pm 0.023$ \\
    low          & 0.004 & $0.958 \pm 0.009$ & $0.087 \pm 0.030$ \\
    medium-low   & 0.020 & $0.952 \pm 0.014$ & $0.103 \pm 0.038$ \\
    medium-high  & 0.100 & $0.934 \pm 0.020$ & $0.153 \pm 0.051$ \\
    high         & 0.500 & $0.708 \pm 0.018$ & $0.198 \pm 0.033$ \\
    \bottomrule
  \end{tabular}
  \caption{LSTM minimal-pair accuracy and attractor gap by training-noise rate
    (mean $\pm$ std over six runs). Accuracy is stable through $10\%$ noise and
    collapses at $50\%$; the attractor gap grows monotonically throughout.}
  \label{tab:main}
\end{table}

\begin{figure}[t]
  \centering
  \includegraphics[width=\linewidth]{acc_vs_noise.png}
  \caption{Overall minimal-pair accuracy vs.\ injected SVA noise rate. The dashed
    line marks chance ($0.5$). Performance is flat through $10\%$ noise and falls
    sharply at $50\%$.}
  \label{fig:acc_vs_noise}
\end{figure}

\subsection{Noise pushes the model toward the linear rule}

Splitting accuracy by attractor presence isolates \emph{which} rule the model
relies on (Figure~\ref{fig:acc_by_attractor}). No-attractor items, where the
hierarchical and linear rules make the same prediction, stay above $0.94$ until
the $50\%$ condition. Attractor items, where the two rules disagree, erode far
faster: accuracy falls from $0.894$ at baseline to $0.525$ (near chance) at
$50\%$ noise. As a result the \emph{attractor gap}, our index of reliance on the
linear, nearest-noun heuristic, grows monotonically with noise, from $0.073$ to
$0.198$ (Figure~\ref{fig:attractor_gap}). Noisier training data drives the model
away from the hierarchical rule and toward the linear shortcut, consistent with
the ``noise pushes models toward simpler generalization'' branch of our second
predicted outcome.

\begin{figure}[t]
  \centering
  \includegraphics[width=\linewidth]{acc_by_attractor.png}
  \caption{Accuracy on attractor vs.\ no-attractor items. The two cells diverge
    as noise rises: attractor accuracy approaches chance while no-attractor
    accuracy stays high until $50\%$.}
  \label{fig:acc_by_attractor}
\end{figure}

\begin{figure}[t]
  \centering
  \includegraphics[width=\linewidth]{attractor_gap.png}
  \caption{Attractor gap (no-attractor minus attractor accuracy) vs.\ noise rate.
    The monotonic rise indicates increasing reliance on the linear cue.}
  \label{fig:attractor_gap}
\end{figure}

\subsection{Learning trajectories: delay vs.\ prevention}

To ask not just whether but \emph{when} the rule is acquired, we checkpoint every
$500$ optimizer steps and re-score the full minimal-pair set at each point
(Figure~\ref{fig:acc_trajectory}). The curves separate the two ways noise can
hurt. For the baseline through medium-high conditions, accuracy follows nearly
the same trajectory and converges to roughly the same endpoint
($0.95$ to $0.97$): noise \emph{delays} acquisition but does not prevent it. The
$50\%$ condition behaves qualitatively differently. It plateaus near $0.69$ and
never reaches the rule, a case of \emph{prevention}. The attractor gap reinforces
this: for the clean model the gap shrinks over training as the hierarchical rule
consolidates ($0.148 \to 0.054$), whereas for the $50\%$ model it \emph{grows}
with more training ($0.143 \to 0.197$), so additional exposure to corrupted data
entrenches the linear heuristic rather than correcting it.

\begin{figure}[t]
  \centering
  \includegraphics[width=\linewidth]{acc_trajectory.png}
  \caption{Minimal-pair accuracy over training steps, one curve per noise
    condition. Low-to-medium noise tracks the clean trajectory and converges
    (delay); the $50\%$ condition plateaus well below the others (prevention).}
  \label{fig:acc_trajectory}
\end{figure}

\subsection{Structure buys robustness: RNNG vs.\ LSTM}
\label{sec:rnng}

To test whether an explicit hierarchical bias changes how a model copes with
noise, we ran the identical five-condition sweep with a Recurrent Neural Network
Grammar (RNNG), which builds an explicit constituency structure as it processes
each sentence. The RNNG's beam-search scorer covers $19{,}983$ of the minimal
pairs against the LSTM's $19{,}819$ (a $\sim$99\% overlap of the same held-out
set), so overall accuracies are comparable up to that small difference in
coverage. The attractor gap, a within-model rate rather than a count, is exact
across architectures and is the cleaner cross-model contrast; an
intersection-matched comparison is left to future work. These RNNG values are a
single full-scale seed (matched to the LSTM protocol at $150$k sentences and
eight epochs). Because seeds are still accruing, we report point estimates
without error bars and read them as a \emph{direction} that the multi-seed run
will tighten.

Two contrasts stand out (Table~\ref{tab:rnng_compare},
Figures~\ref{fig:compare_acc} and \ref{fig:compare_gap}). First, the RNNG is more
accurate than the LSTM at \emph{every} noise level, and its margin widens
exactly where noise bites hardest. The two architectures are nearly tied on
clean data ($0.979$ vs.\ $0.961$), but at $50\%$ noise the RNNG holds $0.834$
where the LSTM has fallen to $0.708$, a $12.6$-point gap. Across the full range
the RNNG sheds only $14.5$ accuracy points from baseline to $50\%$ noise
($0.979 \to 0.834$), against the LSTM's $25.3$-point collapse
($0.961 \to 0.708$). The structural prior does not make the model immune, since
the threshold-like drop at $50\%$ is still visible, but it substantially blunts
it.

Second, and more tellingly, the RNNG barely drifts toward the linear heuristic.
Its attractor gap stays low and essentially flat as noise rises, from $0.020$ at
baseline to $0.036$ at $50\%$ noise (a $1.8\times$ range), whereas the LSTM's gap
grows $2.7\times$, from $0.073$ to $0.198$ (Figure~\ref{fig:compare_gap}). The
contrast is sharp at the extremes: the RNNG's gap under the heaviest noise
($0.036$) is only \emph{half} the LSTM's gap on perfectly clean data ($0.073$).
Where the LSTM trades the hierarchical rule for the nearest-noun shortcut as its
training data degrades, the RNNG keeps resolving agreement structurally almost
regardless of noise. This is the signal our architectural hypothesis predicted:
a built-in syntactic bias preserves the hierarchical generalization under input
corruption that drives a sequential model to the linear one.

\begin{table}[t]
  \centering
  \small
  \begin{tabular}{lccccc}
    \toprule
    & & \multicolumn{2}{c}{Acc.\ (all)} & \multicolumn{2}{c}{Attractor gap} \\
    \cmidrule(lr){3-4}\cmidrule(lr){5-6}
    Condition & Rate & LSTM & RNNG & LSTM & RNNG \\
    \midrule
    baseline    & 0.000 & $0.961$ & $\mathbf{0.979}$ & $0.073$ & $\mathbf{0.020}$ \\
    low         & 0.004 & $0.958$ & $\mathbf{0.979}$ & $0.087$ & $\mathbf{0.030}$ \\
    medium-low  & 0.020 & $0.952$ & $\mathbf{0.978}$ & $0.103$ & $\mathbf{0.022}$ \\
    medium-high & 0.100 & $0.934$ & $\mathbf{0.971}$ & $0.153$ & $\mathbf{0.034}$ \\
    high        & 0.500 & $0.708$ & $\mathbf{0.834}$ & $0.198$ & $\mathbf{0.036}$ \\
    \bottomrule
  \end{tabular}
  \caption{LSTM vs.\ RNNG minimal-pair accuracy and attractor gap by
    training-noise rate, on near-identical held-out sets ($\sim$99\% overlap).
    LSTM values are the mean over
    six runs (Table~\ref{tab:main}); RNNG values are a single full-scale seed.
    The RNNG is more accurate at every noise level and keeps a near-flat,
    low attractor gap where the LSTM's grows steeply.}
  \label{tab:rnng_compare}
\end{table}

\begin{figure}[t]
  \centering
  \includegraphics[width=\linewidth]{combined_acc_vs_noise.png}
  \caption{Overall minimal-pair accuracy vs.\ noise rate for the LSTM and RNNG.
    The RNNG sits above the LSTM throughout and separates sharply at $50\%$
    noise (RNNG single seed; LSTM mean of six).}
  \label{fig:compare_acc}
\end{figure}

\begin{figure}[t]
  \centering
  \includegraphics[width=\linewidth]{combined_attractor_gap.png}
  \caption{Attractor gap vs.\ noise rate. The LSTM's reliance on the linear cue
    climbs steeply with noise ($0.073 \to 0.198$); the RNNG's stays low and
    nearly flat ($0.020 \to 0.036$), below the LSTM's clean-data gap even at
    $50\%$ noise.}
  \label{fig:compare_gap}
\end{figure}
